\section{GAN：隐式密度的经典答案}

\subsection{GAN 为什么重要}

GAN 的出现改变了生成模型的主流叙事。它不再强求写出 $p(x)$，而是让一个生成器 $G(z)$ 从简单噪声出发直接产出样本，再让判别器 $D(x)$ 去判断“真还是假的”。这种方法把生成任务改写成一个博弈问题，直观、简洁，而且在图像上能得到非常锐利的结果。

\begin{figure}[H]
\centering
\includegraphics[width=0.95\textwidth]{slides-images/slide-010.jpg}
\caption{GAN 章节开场：从上一讲的 autoregressive / VAE 自然过渡到隐式密度模型\protect\footnotemark}
\end{figure}
\footnotetext{Slides 第11页。}

\begin{figure}[H]
\centering
\includegraphics[width=0.95\textwidth]{slides-images/slide-013.jpg}
\caption{GAN 的基本数据流：从简单先验 $p(z)$ 采样，经过生成器得到假样本，再交给判别器做真伪判断\protect\footnotemark}
\end{figure}
\footnotetext{Slides 第14页。}

\begin{knowledgebox}{GAN 的关键设定}
\begin{itemize}
    \item $z \sim p(z)$ 是简单先验，通常是标准高斯。
    \item $G(z)$ 把噪声映射成图像或其他样本。
    \item $D(x)$ 输出“这个样本是真实数据的概率”。
    \item 训练目标不是拟合标签，而是让生成分布 $p_G$ 尽量接近真实分布 $p_{\text{data}}$。
\end{itemize}
\end{knowledgebox}

\begin{importantbox}{GAN 和判别式任务的关系}
判别器 $D$ 看起来像一个二分类器，但它不是我们最终要交付的模型。它的角色更像“训练时的评审员”：通过判断真伪来给生成器提供梯度。真正要学的是生成分布，而不是一个分类标签。
\end{importantbox}

\subsection{训练目标不是普通 loss}

GAN 的经典目标是一个 minimax game：

\[
\min_G \max_D \;
\mathbb{E}_{x \sim p_{\text{data}}}[\log D(x)]
\;+\;
\mathbb{E}_{z \sim p(z)}[\log(1 - D(G(z)))]
\]

这个式子很像一个普通 loss，但本质上它不是。它是两个玩家在互相追逐：判别器想把真样本打高分、假样本打低分；生成器想骗过判别器，让假样本看起来像真样本。

\begin{figure}[H]
\centering
\includegraphics[width=0.95\textwidth]{slides-images/slide-016.jpg}
\caption{GAN 的经典 minimax 目标：一个生成器和一个判别器在对抗中共同优化\protect\footnotemark}
\end{figure}
\footnotetext{Slides 第16页。}

\begin{figure}[H]
\centering
\includegraphics[width=0.95\textwidth]{slides-images/slide-021.jpg}
\caption{交替更新伪代码：GAN 不是最小化一个单一标量 loss，而是轮流更新 $D$ 和 $G$\protect\footnotemark}
\end{figure}
\footnotetext{Slides 第21页。}

\begin{warningbox}{为什么 GAN 的训练曲线不好看}
GAN 没有一个稳定、可解释、可单调下降的总 loss。你看到的 $D$ loss 或 $G$ loss 只能说明当前博弈状态的一部分，不能像监督学习那样直接解释为“训练变好了多少”。
\end{warningbox}

\subsection{为什么早期训练会梯度消失}

GAN 初期最常见的情况是：生成器还很弱，判别器一下子就能把真伪分开。于是 $D(G(z)) \approx 0$，这时候如果生成器直接最小化 $\log(1-D(G(z)))$，梯度会很弱，更新信号几乎传不回来。

\begin{figure}[H]
\centering
\includegraphics[width=0.95\textwidth]{slides-images/slide-023.jpg}
\caption{GAN 训练早期的典型问题：判别器过强时，生成器几乎拿不到有用梯度\protect\footnotemark}
\end{figure}
\footnotetext{Slides 第24页。}

\begin{importantbox}{非饱和损失是实战补丁}
实践中常用的 generator loss 不是原始的 $\log(1-D(G(z)))$，而是
\[
\mathcal{L}_G = -\log D(G(z)).
\]
这样做的目的很朴素：当判别器很强时，生成器仍然能收到比较有用的梯度，不至于一开始就“学不动”。
\end{importantbox}

\subsection{为什么这个目标理论上是合理的}

如果判别器容量足够、优化足够理想，那么固定生成器时的最优判别器是

\[
D^*(x)=\frac{p_{\text{data}}(x)}{p_{\text{data}}(x)+p_G(x)}.
\]

把这个 $D^*$ 代回去，外层目标会变成一个和 Jensen-Shannon divergence 相关的量。于是理想情况下，最小化这个博弈等价于让 $p_G$ 逼近 $p_{\text{data}}$。

\begin{figure}[H]
\centering
\includegraphics[width=0.95\textwidth]{slides-images/slide-027.jpg}
\caption{GAN 理论直觉：固定 $G$ 时，最优判别器可以写成一个密度比；理想条件下，外层目标会推动 $p_G$ 向 $p_{\text{data}}$ 靠拢\protect\footnotemark}
\end{figure}
\footnotetext{Slides 第28页。}

\begin{warningbox}{理论正确不等于训练可行}
这个结论依赖很强的理想化前提：无限模型容量、完美优化、有限样本误差可忽略。现实里这些条件都不成立，所以“理论上能证明”并不等于“训练一定顺利”。
\end{warningbox}

\begin{knowledgebox}{GAN 的优点和缺点}
\begin{itemize}
    \item 优点：采样快，视觉效果往往很好，生成结果锐利。
    \item 缺点：训练不稳定，模式塌缩明显，loss 很难解释。
    \item 工程限制：没有显式密度，评估和调参比显式模型更靠经验。
\end{itemize}
\end{knowledgebox}

\subsection{经典架构：DCGAN 和 StyleGAN}

GAN 真正能从 toy example 走向真实图像数据，是靠架构上的一系列经验修正。DCGAN 把卷积引入生成器和判别器，让网络更适合处理图像结构；StyleGAN 则进一步把“风格”注入每一层，显著提升了可控性和视觉质量。

\begin{figure}[H]
\centering
\includegraphics[width=0.95\textwidth]{slides-images/slide-029.jpg}
\caption{DCGAN：GAN 第一次在非 toy 数据上真正跑通的重要架构\protect\footnotemark}
\end{figure}
\footnotetext{Slides 第30页。}

\begin{figure}[H]
\centering
\includegraphics[width=0.95\textwidth]{slides-images/slide-031.jpg}
\caption{StyleGAN 的核心思路：通过自适应归一化把 style 分量注入各层特征\protect\footnotemark}
\end{figure}
\footnotetext{Slides 第32页。}

StyleGAN 里最值得读懂的地方不是“名字叫 style”，而是它对每层特征做了可控的 scale 和 shift。可以把它理解成：生成器不只是拿一个 latent 一次性吐图，而是让 latent 在不同尺度上持续影响图像的纹理、局部结构和语义细节。

\begin{equation*}
\mathrm{AdaIN}(x,w,b)_i = w_i \frac{x_i - \mu(x)}{\sigma(x)} + b_i
\end{equation*}

\begin{knowledgebox}{StyleGAN 的工程含义}
分层注入 style 有两个好处：
\begin{itemize}
    \item 低层更容易控制粗结构，高层更容易控制纹理和细节。
    \item latent space 往往更平滑，便于插值、编辑和风格混合。
\end{itemize}
\end{knowledgebox}

\subsection{Latent interpolation 和生成空间}

GAN 很经典的一个现象是 latent interpolation 平滑。给定两个 latent 向量 $z_0$ 和 $z_1$，如果沿线性路径做插值，再送入生成器，图像往往会连续变形，而不是突然跳变。这说明生成器学到的 latent geometry 至少在局部上是有意义的。

\begin{figure}[H]
\centering
\includegraphics[width=0.95\textwidth]{slides-images/slide-033.jpg}
\caption{GAN 的 latent interpolation：在潜空间中平滑插值，输出图像也通常会平滑过渡\protect\footnotemark}
\end{figure}
\footnotetext{Slides 第34页。}

\begin{importantbox}{这件事为什么重要}
平滑插值不是花哨演示，而是说明生成模型学到的 latent 表示不是随机噪声。它至少在局部上组织了语义结构，所以 latent editing、style mixing 和条件控制才有可能成立。
\end{importantbox}

\subsection{GAN 这一节的结论}

\begin{figure}[H]
\centering
\includegraphics[width=0.95\textwidth]{slides-images/slide-034.jpg}
\caption{GAN 小结：简单目标、很好的样本质量、但训练不稳定且难以扩展\protect\footnotemark}
\end{figure}
\footnotetext{Slides 第35页。}

\begin{importantbox}{GAN 留下的遗产}
GAN 没有消失，它只是退回到了“局部模块”的位置。现代系统仍然会从 GAN 借用：
\begin{itemize}
    \item 对抗损失来提升局部真实感
    \item latent space 里的插值和编辑思路
    \item 对高频细节的偏好
\end{itemize}
\end{importantbox}

